\chapter{Preconditioners}
\label{ch:preconditioners}
\impl{src/linear\_solvers.jl}

A preconditioner is an approximation $\Prec \approx \Kmat^{-1}$ that is cheap
to apply and reduces $\kappa(\Prec\Kmat)$ in \eqref{eq:cg-bound}.  It must be
symmetric positive definite for conjugate gradients to remain valid.  It need
not be accurate: Remark~\ref{ch:constitutive} applies here too --- the
preconditioner shapes the path, never the fixed point.

\section{Jacobi}

$\Prec = \Dmat^{-1}$ with $\Dmat = \diag(\Kmat)$.  One vector scaling per
application, trivially parallel, always available.  It removes the
conditioning that comes from heterogeneous element sizes and material
properties, and nothing else --- in particular it does nothing about the
$(L/h)^2$ growth with mesh refinement that dominates quasi-static problems.

Its unusual effectiveness in implicit dynamics at small $\Delta t$ is a direct
consequence of \eqref{eq:cM}: when $c_M\Mmat$ dominates $\Keff$, and $\Mmat$ is
nearly diagonal for the element types in use, the diagonal captures most of the
operator.

\section{Incomplete \texorpdfstring{$LDL^{\mathsf{T}}$}{LDLT}}

An incomplete factorization $\Kmat \approx \tensor{L}\Dmat\tensor{L}^{\mathsf{T}}$
that drops fill outside a chosen sparsity pattern, applied by two triangular
solves.  Substantially stronger than Jacobi for ill-conditioned systems, and
the natural choice for $J_2$ plasticity on non-uniform meshes.

It is CPU-only, and the obstruction is structural rather than a missing port: a
triangular solve is inherently sequential along each dependency chain, which is
the one access pattern a GPU cannot exploit.  This is why the GPU path needs
either a polynomial preconditioner or multigrid, both of which reduce to
matrix--vector products.

\section{Chebyshev polynomial preconditioner}
\label{sec:chebyshev}

A polynomial preconditioner approximates $\Kmat^{-1}$ by a polynomial in
$\Kmat$, so applying it costs only matrix--vector products --- which makes it
GPU-friendly where triangular solves are not.

Carina implements Chebyshev polynomials of the \emph{fourth kind} with optimal
weights, following Lottes~\cite{Lottes2022} as implemented in Ifpack2.  This
variant has three properties that matter here:

\begin{itemize}
  \item It requires only $\lambda_{\max}$, not the eigenvalue \emph{ratio}.
        Standard Chebyshev needs a lower bound $\lambda_{\min}$ as well, and a
        poor estimate of it degrades the polynomial badly; $\lambda_{\min}$ is
        also the harder of the two to estimate.
  \item Jacobi scaling is baked into the recurrence, so $\Dmat^{-1}$ is applied
        at each stage rather than as a separate wrapper.
  \item The optimal weights make the resulting operator symmetric positive
        definite \emph{by construction}, so it is admissible as a CG
        preconditioner without the squaring trick --- costing $k$ matrix--vector
        products per application instead of $2k$.
\end{itemize}

With $\sigma = 1/\lambda_{\max}$ and weights $\beta_i$ tabulated for degrees
$1$ through $16$:

\begin{algorithm}[h]
\caption{Fourth-kind Chebyshev preconditioner application, degree $k$}
\label{alg:chebyshev}
\begin{algorithmic}[1]
  \Require $\vect{b}$, operator $\Kmat$, $\Dmat^{-1}$, $\sigma = 1/\lambda_{\max}$,
           weights $\beta_1,\dots,\beta_k$
  \Ensure $\vect{y} \approx \Kmat^{-1}\vect{b}$
  \State $\vect{z} \gets \tfrac43\sigma\,\Dmat^{-1}\vect{b}$,\quad
         $\vect{x} \gets \vect{z}$,\quad
         $\vect{y} \gets \beta_1\vect{z}$
  \For{$i = 1$ \textbf{to} $k-1$}
    \State $\gamma_i \gets \dfrac{2i-1}{2i+3}$,\qquad
           $\rho_i \gets \dfrac{8i+4}{2i+3}\,\sigma$
    \State $\vect{z} \gets \rho_i\,\Dmat^{-1}\left(\vect{b} - \Kmat\vect{x}\right)
                            + \gamma_i\vect{z}$
      \Comment{one operator application}
    \State $\vect{x} \gets \vect{x} + \vect{z}$
      \Comment{raw fourth-kind iterate}
    \State $\vect{y} \gets \vect{y} + \beta_{i+1}\vect{z}$
      \Comment{weighted accumulation}
  \EndFor
  \State \Return $\vect{y}$
\end{algorithmic}
\end{algorithm}

\noindent Note that $\vect{x}$ and $\vect{y}$ are distinct: $\vect{x}$ carries
the unweighted recurrence that generates the polynomial basis, while $\vect{y}$
accumulates the optimally weighted combination that is actually returned.
Collapsing them --- returning $\vect{x}$ --- gives the ordinary fourth-kind
polynomial without the optimal weights, and loses the positive-definiteness
that makes the operator admissible to conjugate gradients.

\paragraph{Spectral bound.}  $\lambda_{\max}$ of $\Dmat^{-1}\Kmat$ is estimated
by ten power iterations followed by a Rayleigh quotient, then inflated by a
safety factor of $1.1$.  The inflation is deliberate and one-sided:
overestimating $\lambda_{\max}$ makes the polynomial conservative and merely
slower, whereas underestimating it puts part of the spectrum outside the
interval the polynomial damps, and the preconditioner amplifies those modes
instead of removing them.

\paragraph{Symmetric scaling.}  On the assembled path the polynomial is built
on $\Dmat^{-1/2}\Kmat\Dmat^{-1/2}$ rather than on $\Kmat$, for the reason given
in Section~\ref{sec:elimination}: penalty-scale entries give $\Kmat$
eigenvalues of order $10^{15}$ while the symmetrically scaled operator has
eigenvalues of order one, and a polynomial fitted to the former is useless.
The preconditioner is then $\Prec = \Dmat^{-1/2} p_k(\tensor{S})\, \Dmat^{-1/2}$,
which is symmetric positive definite whenever $p_k$ is.
